\subsection{Matched response under alternative crossings, support, and weighting}
\label{sec:v20-mrc}
We use the same background feature-map LPIPS outcome in both studies. Let $d_i$ be the signed area between two preservation curves divided by their positive response-overlap width $w_i$ on source $i$. The primary estimate is the mean of $d_i$. Weighting by width gives $\sum_i w_i d_i/\sum_i w_i$, so a source spanning twice the response range receives twice the weight. Positive differences indicate lower background LPIPS for the subtracted arm.

\begin{table*}[t]
\centering\footnotesize
\caption{SAM matched-response sensitivity. Each estimate is a difference in background feature-map LPIPS with its marginal 95\% paired-source bootstrap interval. S, U and F denote spatial SAM, uniform SAM and FlowEdit. For the last two rows, $n$ counts sources contributing to the mean curves.}
\label{tab:v20-mrc}
\setlength{\tabcolsep}{3pt}
\begin{tabular}{@{}lrlrlrl@{}}
\hline
 & \multicolumn{2}{c}{U $-$ S} & \multicolumn{2}{c}{F $-$ S} & \multicolumn{2}{c}{F $-$ U}\\
Analysis & $n$ & Mean [95\% CI] & $n$ & Mean [95\% CI] & $n$ & Mean [95\% CI]\\
\hline
First crossing & 41 & $0.0390\;[0.0106,0.0632]$ & 38 & $0.1196\;[0.0739,0.1631]$ & 41 & $0.0638\;[0.0229,0.1046]$\\
Last crossing & 41 & $0.0457\;[0.0134,0.0729]$ & 38 & $0.1428\;[0.0919,0.1914]$ & 41 & $0.0761\;[0.0344,0.1196]$\\
Width weighted & 41 & $0.0696\;[0.0435,0.0946]$ & 38 & $0.1419\;[0.1000,0.1862]$ & 41 & $0.0663\;[0.0162,0.1159]$\\
Common sources & 38 & $0.0429\;[0.0124,0.0682]$ & 38 & $0.1196\;[0.0739,0.1631]$ & 38 & $0.0589\;[0.0167,0.1028]$\\
Common interval & 38 & $0.0467\;[0.0079,0.0797]$ & 38 & $0.1162\;[0.0700,0.1609]$ & 38 & $0.0695\;[0.0323,0.1084]$\\
Common interval, width weighted & 38 & $0.0707\;[0.0432,0.0981]$ & 38 & $0.1441\;[0.0970,0.1937]$ & 38 & $0.0734\;[0.0274,0.1176]$\\
Two prompts excluded & 39 & $0.0385\;[0.0087,0.0639]$ & 36 & $0.1161\;[0.0673,0.1621]$ & 39 & $0.0591\;[0.0175,0.1022]$\\
Edit MAE response & 42 & $0.0417\;[0.0227,0.0598]$ & 41 & $0.0834\;[0.0584,0.1092]$ & 42 & $0.0260\;[0.0089,0.0436]$\\
Mean curves, pair sources & 41 & $0.0758\;[0.0424,0.1006]$ & 38 & $0.1194\;[0.0722,0.1837]$ & 41 & $0.0387\;[-0.0020,0.0926]$\\
Mean curves, all eligible & 42 & $0.0744\;[0.0403,0.0993]$ & 42 & $0.1190\;[0.0732,0.1828]$ & 42 & $0.0426\;[0.0017,0.0990]$\\
\hline
\end{tabular}
\end{table*}

\paragraph{Coverage and integration.}
On SAM, 42 of 50 sources have a background. At first directional-CLIP crossing, 41, 38 and 41 have positive overlap for U$-$S, F$-$S and F$-$U, respectively. The median widths are 0.07594, 0.09204 and 0.09399 directional-CLIP units. In each case the remaining eligible sources have disjoint ranges. We connect the five observations in setting order, split the overlap at response knots, and integrate each linear piece. First and last crossing select the earliest and latest segments covering its midpoint. For the ambiguity check we omit sources \texttt{611000000001} and \texttt{812000000002}, leaving 40 background-eligible sources among 48. The MAE rows match the magnitude of RGB change within the edit mask.

\paragraph{Source set, interval and weighting.}
All three comparisons share 38 sources. On these sources, U$-$S is $+0.0429$ over each pair's own interval and $+0.0467$ over each source's three-arm interval. Changing only the interval shifts the estimate by $+0.0039$ $[-0.0087,0.0206]$. The three-arm intervals have median width 0.05298. Integrating every pair over those intervals gives additive contrasts within each source and in the means. On the primary 41 sources, weighting by pairwise width increases U$-$S by $+0.0306$ $[0.0129,0.0493]$; taking the last crossing changes it by $+0.0067$ $[-0.0023,0.0174]$. Computing the two cohort-mean curves first on those same 41 sources increases U$-$S by $+0.0368$ $[0.0062,0.0642]$.

\paragraph{Restoring operator.}
We first average seeds 0 and 17 at each of three measured scales (0.7, 1.0 and 1.3) for both allocations on the uniform-origin profile. At equal directional CLIP, 36 of 40 sources have positive overlap (median width 0.01286). Uniform minus request LPIPS is $+0.0128$ $[0.0022,0.0240]$ at first crossing, $+0.0178$ $[0.0061,0.0304]$ at last crossing and $+0.0121$ $[0.0012,0.0243]$ under width weighting. Matching edit MAE gives $+0.0308$ $[0.0197,0.0424]$ on 40 sources (median width 0.03184 MAE). Averaging the CLIP curves before matching gives disjoint ranges on all 40. On the 36 individually supported sources the mean curves overlap by 0.00032 CLIP, yielding $+0.0117$; 7,792 of 10,000 resamples support that estimate. The historical inverse comparison matches mean curves on background LPIPS and measures directional CLIP. Its 40-source uniform minus request CLIP difference is $+0.00652$ $[0.00086,0.01266]$ over LPIPS 0.43297 to 0.53186.

\paragraph{Uncertainty.}
The forward comparisons use 10,000 paired source resamples and percentile intervals. Each draw carries all settings and arms for a sampled source. We recompute width-weighted ratios and mean-curve overlaps in each draw. We report an interval when all draws support the estimate. For differences between estimators with different contributing sets, we resample their union and apply each estimator's inclusion rule to the same draw. The accompanying CSV also gives contrasts restricted to the source intersection. The historical inverse comparison retains its original 20,000 category-stratified source draws.
